\section{基于KG-GCN的燃煤机组汽轮机系统故障诊断方法}
\label{sec:kg_gcn_v5}

\subsection{引言}

汽轮机系统发生故障时，异常状态会沿主通流、再热、抽汽回热及冷端等实际热力路径向相邻设备传播。若故障未能及时识别，不仅可能造成局部设备性能恶化，还可能进一步影响机组安全性和经济性。现有故障诊断方法大体可分为两类：一类是基于运行数据的故障分类方法，如支持向量机、卷积神经网络和自编码器等，该类方法能够从历史样本中自动提取故障特征，但模型性能高度依赖故障样本的数量、代表性和标签质量；另一类是基于机理或专家经验的规则诊断方法，该类方法具有较强的物理解释能力，但面对复杂变工况和多参数耦合故障时，固定规则难以完整覆盖实际运行状态。对于燃煤机组汽轮机系统而言，高质量现场故障样本数量有限，而大量设备结构、测点归属、运行检修知识和故障机理又长期分散存在于设备资料、DCS点表、检修记录以及公开技术文献中，因此，仅依赖运行数据或仅依赖经验规则均难以充分利用现场已经积累的多源信息。

第三章将汽轮机设备拓扑、DCS参数映射以及故障机理知识组织为结构化知识图谱，使设备、参数、异常和故障之间的静态关联能够以图结构形式表示；第二章建立的健康状态模型则能够给出当前运行边界下各状态参数的健康参考，并形成实测状态与健康模型输出之间的动态偏差信息。基于此，本文在课题组既有KG-GCN技术路线\cite{YangHaodong2026CFB}基础上，将运行数据与静态领域知识映射到同一参数空间，构建面向燃煤机组汽轮机系统的知识图谱图卷积网络（Knowledge Graph--Graph Convolutional Network，KG-GCN）故障诊断方法。该方法首先由知识图谱生成参数关联邻接矩阵，再将状态监测残差窗口加载至图谱参数节点，通过图卷积逐层聚合节点及其邻域特征；针对高质量故障标签稀缺的问题，进一步采用残差掩码预训练学习参数之间的结构关联，并在预训练结束后固定图卷积编码器，利用带标签故障窗口完成故障分类。

\subsection{图卷积网络}
\label{subsec:gcn_theory_v5}

复杂工业热力系统的参数关联本质上属于非欧几里得拓扑结构。主蒸汽、再热、抽汽和凝汽器冷端参数之间的联系由实际设备连接和介质传递路径决定，并不存在类似规则图像网格的固定空间邻接关系。传统卷积神经网络适合处理欧几里得空间中的规则局部结构，而图卷积网络（Graph Convolutional Network，GCN）能够直接在图拓扑上聚合节点自身及其空间邻域信息，从而捕获参数之间的结构化依赖关系\cite{Kipf2017GCN}。

设给定参数图$G=(V,E)$，其中$V=\{v_1,v_2,\ldots,v_N\}$为参数节点集合，$E$为节点之间的关联边集合。图的拓扑连接关系由邻接矩阵$A\in\mathbb{R}^{N\times N}$表示，节点基础属性由特征矩阵$H^{(0)}$表示。在图卷积计算过程中，除相邻节点信息外，节点自身状态也应得到保留，因此在邻接矩阵中加入单位矩阵$I_N$，得到增强邻接矩阵
\begin{equation}
\tilde A=A+I_N .
\end{equation}
相应的度矩阵定义为
\begin{equation}
\tilde D_{ii}=\sum_j\tilde A_{ij} .
\end{equation}
对增强邻接矩阵进行对称归一化可得
\begin{equation}
\bar A=\tilde D^{-\frac{1}{2}}\tilde A\tilde D^{-\frac{1}{2}} .
\end{equation}
标准图卷积层的特征传播过程可写为
\begin{equation}
H^{(l+1)}=\sigma\left(\bar A H^{(l)}W^{(l)}\right),
\end{equation}
式中：$H^{(l)}$为第$l$层节点特征；$W^{(l)}$为第$l$层可学习权重矩阵；$\sigma(\cdot)$为非线性激活函数。由上述计算过程可知，邻接矩阵决定特征传播路径，权重矩阵完成节点特征维度变换，每经过一层图卷积，目标节点均能够融合自身与一阶邻域信息；随着图卷积层数增加，节点可以进一步获得更远邻域中的关联参数信息。

第五章VHP代表性算例采用两层图卷积编码器。首先将每个参数节点的连续残差窗口投影至$d$维特征空间
\begin{equation}
H^{(0)}=\mathrm{ReLU}(R_tW_e+b_e),
\end{equation}
随后利用两层GCN进行参数关联信息聚合。考虑固定DCS节点自身残差具有明确物理含义，为减小多层传播对节点原始信息的削弱，本文在图卷积层中加入残差连接和LayerNorm：
\begin{equation}
H^{(1)}=\mathrm{ReLU}\left[\mathrm{LN}\left(H^{(0)}+\bar A H^{(0)}W^{(1)}\right)\right],
\end{equation}
\begin{equation}
H^{(2)}=\mathrm{ReLU}\left[\mathrm{LN}\left(H^{(1)}+\bar A H^{(1)}W^{(2)}\right)\right].
\end{equation}
第一层主要融合直接相邻参数的信息，第二层进一步聚合两跳邻域特征，使节点表示同时包含自身动态状态与局部热力拓扑信息。

% 图4-x：GCN框架示意图。最终绘图仅根据审计后的邻接关系生成。

\subsection{KG-GCN模型输入构建}
\label{subsec:kg_input_v5}

第三章构建的汽轮机系统知识图谱同时包含设备、参数、异常、故障和处置等多类实体，而实时图卷积输入需要与DCS数值建立一一对应关系。因此，首先从系统级知识图谱中提取能够映射至实时状态量的参数节点，并根据设备归属、主通流接口、抽汽支路及相邻热力边界等关系构造参数关联图。设用于故障诊断的参数图为
\begin{equation}
G_p=(V_p,E_p),
\end{equation}
式中：$V_p$为能够与DCS状态参数一一对应的节点集合；$E_p$为由知识图谱关系投影得到的参数关联边集合。由$E_p$可进一步得到参数邻接矩阵$A_p$，并按照前述方法得到归一化邻接矩阵$\bar A_p$。该矩阵反映运行生产资料中保存的参数关联结构，是KG-GCN模型的静态结构输入。

在第二章健康状态模型得到各参数预测结果后，单一时刻的状态残差能够表示当前实际运行状态相对于健康参考状态的瞬时偏离。对于第$i$个状态参数，其在时刻$t$的实测值和健康模型输出分别记为$y_i(t)$和$\hat y_i(t)$，则状态残差定义为
\begin{equation}
e_i(t)=y_i(t)-\hat y_i(t) .
\end{equation}
不同参数量纲和波动范围存在差异，因此采用训练集健康残差统计量进行标准化：
\begin{equation}
r_i(t)=\frac{e_i(t)-\mu_i}{\sigma_i+\varepsilon},
\end{equation}
式中：$\mu_i$和$\sigma_i$分别为训练集第$i$个参数健康残差的均值和标准差；$\varepsilon$为防止分母为零的微小常数。验证集和测试集不参与标准化参数估计。

设参与故障诊断的参数数量为$N$，在时刻$t$得到的标准化残差向量表示为
\begin{equation}
\mathbf r_t=[r_1(t),r_2(t),\ldots,r_N(t)]^{\mathrm T} .
\end{equation}
该向量能够表征各参数在当前时刻相对于健康状态的偏离程度，但仅使用单一时刻残差容易忽略故障发生和演化过程中的连续变化。为此，采用滑动时间窗口整合一段时间内的残差序列。设窗口长度为$L_r$，则在时刻$t$结束的残差窗口矩阵为
\begin{equation}
R_t=
\begin{bmatrix}
r_1(t-L_r+1)&\cdots&r_1(t)\\
r_2(t-L_r+1)&\cdots&r_2(t)\\
\vdots&\ddots&\vdots\\
r_N(t-L_r+1)&\cdots&r_N(t)
\end{bmatrix}
\in\mathbb{R}^{N\times L_r} .
\end{equation}
矩阵的行与参数图中的节点一一对应，列表示窗口内的时间位置。由此，归一化邻接矩阵$\bar A_p$作为图结构输入，残差窗口矩阵$R_t$作为节点动态特征输入，二者共同构成KG-GCN的基本输入形式，实现第三章静态知识关系与第二章动态状态监测结果在统一参数空间中的融合。第五章VHP算例中，$N=24$，$L_r=60$。

% 图4-x：KG-GCN模型输入构建流程，应对应“KG邻接矩阵 + 残差窗口矩阵”两路输入。

\subsection{KG-GCN模型训练}
\label{subsec:kg_training_v5}

实际电站运行过程中，正常和一般变工况数据能够长期连续积累，而能够清晰标注故障类型和故障时段的高质量样本数量有限。若直接利用少量故障标签训练图卷积网络，模型容易过度拟合有限故障模式。为充分利用大量未标注状态残差，本文采用“无标签残差掩码预训练--固定图卷积编码器--带标签故障分类”的两阶段训练方式\cite{YangHaodong2026CFB}。

预训练阶段不使用故障类别标签。设无标签残差窗口样本集合为
\begin{equation}
\mathcal D_u=\{R^{(1)},R^{(2)},\ldots,R^{(M_u)}\},
\end{equation}
其中$M_u$为无标签残差窗口数量。对于任一残差窗口$R$，随机生成掩码矩阵$M\in\{0,1\}^{N\times L_r}$，若$M_{ij}=0$，则对应残差位置被遮蔽；若$M_{ij}=1$，则该位置保持可见。掩码输入表示为
\begin{equation}
\tilde R=M\odot R+(1-M)\odot R_{mask},
\end{equation}
式中：$\odot$为逐元素乘法；$R_{mask}$为掩码替代值，本文取0。第五章算例随机遮蔽比例为15\%。

掩码残差窗口和归一化邻接矩阵共同输入图卷积编码器
\begin{equation}
Z=E_{\theta}(\bar A_p,\tilde R),
\end{equation}
并通过解码器恢复原始残差窗口
\begin{equation}
\hat R=D_{\phi}(Z) .
\end{equation}
为了使模型重点学习被遮蔽位置与其邻域之间的关联，预训练损失只在被遮蔽位置计算：
\begin{equation}
\mathcal L_{mask}=\frac{\sum_{ij}(1-M_{ij})(R_{ij}-\hat R_{ij})^2}{\sum_{ij}(1-M_{ij})+\varepsilon} .
\end{equation}
通过最小化掩码重构误差，图卷积编码器需要利用未遮蔽参数及其知识图谱邻域恢复缺失信息，从而在不依赖故障标签的条件下学习多参数之间的结构化残差关系。

预训练完成后固定图卷积编码器参数，仅利用带标签故障窗口训练分类层。设有标签样本集合为
\begin{equation}
\mathcal D_l=\{(R^{(n)},y^{(n)})\}_{n=1}^{M_l},
\end{equation}
其中$y^{(n)}$为第$n$个样本的故障类别标签。为保留固定DCS节点的物理身份，本文不对全部参数节点进行全局平均池化，而将节点局部表示$H^{(0)}$与两层图传播结果$H^{(2)}$按固定节点顺序组合并展开为图级特征向量
\begin{equation}
\mathbf g=\mathrm{vec}\left(\mathrm{LN}[H^{(0)}\Vert H^{(2)}]\right) .
\end{equation}
随后通过全连接层和Softmax得到故障类别概率
\begin{equation}
\hat{\mathbf p}=\mathrm{Softmax}\left(W_c\,\mathrm{ReLU}(W_r\mathbf g+b_r)+b_c\right) .
\end{equation}
分类阶段采用交叉熵损失函数
\begin{equation}
\mathcal L_{cls}=-\frac{1}{M_l}\sum_{n=1}^{M_l}\sum_{k=1}^{C}y_{nk}\log \hat p_{nk} .
\end{equation}
训练结束后，将实时状态监测残差按照相同方式构造为滑动窗口，并与知识图谱邻接矩阵共同输入已训练KG-GCN，最终根据Softmax输出完成故障类别识别。

% 图4-x：KG-GCN训练流程，应突出“Mask预训练”和“固定GCN后分类”两个阶段。

\subsection{本章小结}

本章提出了基于KG-GCN的燃煤机组汽轮机系统故障诊断方法。首先介绍图卷积网络的节点特征传播过程，并利用第三章知识图谱中的设备与参数关联关系构造参数邻接矩阵；随后将第二章健康状态模型产生的标准化残差按滑动窗口形式映射至图谱参数节点，使静态知识关系和动态运行状态在统一图结构中实现融合；针对现场故障标签有限的问题，进一步采用残差掩码重构进行无标签预训练，预训练完成后固定图卷积编码器，并利用带标签故障样本训练分类层。下一章将选取VHP作为代表性设备，对CNN、AE和KG-GCN的故障诊断结果进行对比验证。